\section{Discussion}
\label{sec:discussion}

\subsection{What the two studies establish}

Study T shows that the contract can keep target, provider, capability, action,
and evidence aligned across the evaluated client--backend boundary.  The
direct-wiring comparison also bounds this claim: under the shared policy, both
representations preserve the same ownership and routing decisions and detect
the same common faults.  FunctionalMLDS adds value where relations cross
artifacts, because provider, capability, binding, action, and assertion links
become typed invariants that can be checked and localized.  It does not by
itself make the routing policy more correct.

Study U adds the complementary interface result.  Participants generally
reported that answers matched selected objects and that an observed handoff and
final agent were clear.  However, only 37/65 identified the intended
object-and-topic responsibility rule, and ratings for the usefulness of several
specialized agents were mixed.  Technical enforcement is therefore necessary
but not sufficient for intelligibility: the interface can execute the correct
provider transition while leaving its rationale only partly understood.

\subsection{When the added structure is justified}

A small, static room with one general agent may need only direct configuration
and tests.  The contract is better motivated when several agents divide
responsibility, scenes are regenerated, one operation has many target-specific
bindings, or failures must be traced across scene and backend artifacts.  This
assurance has a cost: the representation stores more references and the
normalized V2 workload is slower than direct wiring in the reported
microbenchmark.  The absolute workload remains small, but it is not end-to-end
latency, and the benchmark does not measure authoring effort.  We therefore
argue for scoped adoption rather than universal replacement of simpler
configurations.

\subsection{Interface implications and next evidence}

The contract already supplies the states needed to communicate the selected
target, responsible provider, handoff reason, and rejection location.  The user
results suggest that the current lightweight cues communicate successful
handoffs more effectively than the underlying division of responsibility.
Near-term improvements should therefore make agent roles and provider changes
more explicit, for example through concise role summaries, a short reason at
handoff time, and clarification choices when a request cannot proceed.

The present user tasks covered successful grounding and one permitted handoff,
not ambiguity, unavailable capabilities, unreachable providers, or rollback.
A controlled failure-comprehension task is the most direct next interface
evaluation.  On the technical side, an independently authored scene with
natural group- and zone-level responsibility, unassigned and ambiguous objects,
and an independently reviewed route oracle would strengthen transfer evidence.
A later developer study could then test whether typed diagnostics reduce fault
localization and repair effort.

\subsection{Deployment boundary}

``Permitted'' means declared by the interaction model; it does not authenticate
a visitor or authorize access to protected resources.  A deployment should
compose the resolved target--provider--action tuple with an external security
policy.  It should also minimize stored utterance and spatial data, provide
retention controls, and review modeled responsibilities for inappropriate or
biased authority.  Consistent assignment is inspectable, but not automatically
fair or safe.
